\documentclass[usegeometry=true]{scrartcl}

% Pakete, Layoutmasse, Listings-Stil und Makros
\usepackage[ngerman]{babel}
\usepackage[T1]{fontenc}
\usepackage{lmodern}
\usepackage[utf8]{inputenc}
\usepackage{hyperref}
\usepackage{amssymb}
\usepackage{amsmath}
% Dimensionen bitte nicht ändern.
\usepackage[left=2cm, right=2cm, top=2cm, bottom=2cm, bindingoffset=1cm, includeheadfoot]{geometry}
%Zeilenabstand bitte nicht ändern
\usepackage[onehalfspacing]{setspace}

\usepackage[backend=biber,style=numeric,]{biblatex}\addbibresource{literatur.bib}

% --- zusätzliche Pakete (Layoutmaße der Vorlage bleiben unverändert) --------
\usepackage{csquotes}
\usepackage{graphicx}
\usepackage{booktabs}
\usepackage{array}
\usepackage{xcolor}
\usepackage{float}
\usepackage{verbatim}
\usepackage{listings}
\usepackage{microtype}
\usepackage[font=small]{caption}

\hypersetup{colorlinks=true, linkcolor=black, citecolor=black, urlcolor=blue!60!black}

% --- Angaben zur Abgabe ----------------------------------------------------
\newcommand{\repourl}{https://github.com/yaniktfl/iv-repo}
\newcommand{\repobranch}{main}
\newcommand{\repocommit}{ce20021}

% --- Elm-Listings ----------------------------------------------------------
\lstdefinelanguage{Elm}{
  morekeywords={module,exposing,import,type,alias,let,in,if,then,else,case,of,port,as,where,True,False},
  morekeywords=[2]{Model,Msg,Maybe,Just,Nothing,List,Set,Dict,Float,Int,String,Bool,Html,Svg,Cmd,Sub,Result,Ok,Err},
  sensitive=true,
  morecomment=[l]{--},
  morecomment=[n]{\{-}{-\}},
  morestring=[b]",
}
\lstset{
  language=Elm,
  basicstyle=\ttfamily\footnotesize,
  keywordstyle=\color{blue!55!black}\bfseries,
  keywordstyle=[2]\color{teal!70!black},
  commentstyle=\color{gray}\itshape,
  stringstyle=\color{purple!60!black},
  showstringspaces=false,
  columns=fullflexible,
  breaklines=true,
  frame=single,
  rulecolor=\color{gray!50},
  framesep=4pt,
  xleftmargin=6pt,
  captionpos=b,
  aboveskip=0.8em,
  belowskip=0.5em,
  literate={ä}{{\"a}}1 {ö}{{\"o}}1 {ü}{{\"u}}1 {Ä}{{\"A}}1 {Ö}{{\"O}}1 {Ü}{{\"U}}1 {ß}{{\ss}}1
           {€}{{\texteuro}}1 {·}{{$\cdot$}}1 {→}{{$\rightarrow$}}1 {Ø}{{\O}}1 {–}{{--}}1,
}

% --- Abbildungsmakro -------------------------------------------------------
% #1 Dateiname ohne Endung in figures/, #2 Breite als Anteil von \textwidth,
% #3 Bildunterschrift, #4 Beschreibung für den Platzhalter.
\newcommand{\abb}[4]{%
  \begin{figure}[H]
    \centering
    \IfFileExists{figures/#1.png}
      {\includegraphics[width=#2\textwidth]{figures/#1}}
      {\fbox{\begin{minipage}[c][0.16\textheight][c]{0.92\textwidth}
        \centering\small
        \textbf{Platzhalter für \texttt{figures/#1.png}}\\[0.5em]
        #4
      \end{minipage}}}
    \caption{#3}\label{fig:#1}
  \end{figure}}


\begin{document}

% ----------------------------------------------------------------------------
\subject{Projektbericht zum Modul Information Visualisierung Sommersemester 2026}
\title{Wann deckt der Wind, wann die Sonne?}
\subtitle{Eine interaktive Visual-Analytics-Anwendung für Erzeugung, Preis und Handel im deutschen Stromsystem 2024}
\author{Lukas Schlinsog (222212780) \and Yanik Engelhardt (218205378)}
\date{\today}
\maketitle
% ----------------------------------------------------------------------------

\begin{center}
\begin{minipage}{0.86\textwidth}
\small
\textbf{Abgabe.} Repository: \url{\repourl} \quad|\quad Branch: \texttt{\repobranch} \quad|\quad
Bewertungsstand: Commit \texttt{\repocommit}. Die Anwendung liegt in \texttt{projekt/}, der Bericht in
\texttt{bericht/}. Die Anwendung übersetzt fehlerfrei mit
\texttt{elm make src/Main.elm -{}-output=main.js} (Elm 0.19.2) und lädt ihre Daten zur Laufzeit per HTTP.
\end{minipage}
\end{center}

% ----------------------------------------------------------------------------
% Inhaltsverzeichnis:
\tableofcontents
\newpage
% ----------------------------------------------------------------------------
% Gliederung und Text:

% ----------------------------------------------------------------------------
% Kapitel -- jedes liegt als eigene Datei in kapitel/ und laesst sich einzeln
% bearbeiten und versionieren.
% ----------------------------------------------------------------------------
\section{Einleitung}

Im deutschen Stromsystem ist der Beitrag der erneuerbaren Energien 2024 der Normalfall geworden.
Im hier verwendeten Datensatz decken sie im Jahresmittel $55{,}1\,\%$ der Last. Dieser Mittelwert
ist zugleich die am wenigsten aussagekräftige Zahl des Datensatzes, denn er entsteht aus Stunden
zwischen $12{,}0\,\%$ und $138{,}6\,\%$, aus $441$ Stunden mit negativen Börsenpreisen und aus
einzelnen Stunden über $900$\,EUR/MWh. Für das Verständnis des Gesamtsystems ist daher die zeitliche
Verteilung der Werte ausschlaggebend.

Darin liegt das Zielproblem dieses Projekts. Auf physikalischer und marktlicher Ebene beschreiben sechs kontinuierliche Zeitreihen den Zustand des Gesamtsystems, namentlich Last, Solarerzeugung, Winderzeugung, Erneuerbarenanteil, Börsenpreis sowie grenzüberschreitender Handel. Für die multidimensionale Betrachtung ganzer Tage tritt als siebtes Attribut die Anzahl negativer Preisstunden hinzu. Jede dieser Größen ist einzeln leicht verfügbar. Schwer zu beobachten ist ihr Zusammenspiel. Dieses entfaltet sich auf drei Zeitskalen gleichzeitig. Dies betrifft den Tagesgang
mit der solaren Mittagsspitze, das mehrtägige Wettergeschehen mit Windphasen und Dunkelflauten sowie
den übergeordneten Jahresgang. Etablierte Darstellungen erzwingen zumeist eine Festlegung auf eine
dieser Skalen. Während eine Jahresübersicht den Tagesgang einebnet, verliert eine isolierte Tagesansicht
den saisonalen Kontext. Die Frage \enquote{Sind negative Preise ein Solarphänomen?} erfordert jedoch beide
Skalen zugleich und zusätzlich einen Weg, ganze Tage miteinander zu vergleichen.

Daraus ergeben sich vier Fragestellungen, die mit Techniken der Informationsvisualisierung
beantwortet werden können und an denen diese Arbeit zu messen ist.

\begin{itemize}
  \item[\textbf{F1}] Wie verteilt sich die erneuerbare Deckung über Jahreszeit und Tageszeit, und
        welche der beiden Skalen erklärt mehr von der Schwankung?
  \item[\textbf{F2}] Wann treten Extremzustände auf? Gemeint sind hohe Deckung mit negativen
        Preisen auf der einen Seite und Dunkelflauten mit hohen Preisen auf der anderen. Sind das
        Einzelstunden oder Tagesmuster?
  \item[\textbf{F3}] Wie tragen Wind- und Solarenergie jeweils dazu bei? Sind sie austauschbar oder
        tragen sie zu verschiedenen Zeiten und mit verschiedenen Folgen bei?
  \item[\textbf{F4}] Wie hängen Erneuerbarenanteil, Preis und Nettohandel an einem konkreten Tag
        zusammen, und wie typisch ist dieser Tag im Vergleich zu allen anderen?
\end{itemize}

Keine dieser Fragen lässt sich mit einer einzelnen Kennzahl oder einem isolierten Diagramm
beantworten. So erfordert F1 zwei Zeitskalen gleichzeitig, F2 das Auffinden seltener Ereignisse in
$8\,690$ Messwerten und F3 sowie F4 den Vergleich mehrerer Attribute mit unterschiedlichen physikalischen
Einheiten. Ein reales Anwendungsproblem, für das Visualisierungstechniken ausgewählt, kombiniert und
begründet werden, entspricht dem Konzept einer \emph{Design Study} nach Munzner \cite{Munzner2008}.

\subsection{Anwendungshintergrund}

Für die Einordnung der Fragestellungen sind vier Zusammenhänge des Anwendungsfeldes maßgeblich.

\textbf{Dargebotsabhängige Erzeugung.} Solar- und Windanlagen speisen abhängig vom Wettergeschehen
ein, losgelöst vom aktuellen Strombedarf. Die Photovoltaik folgt einem deterministischen Tages- und
Jahresgang mit dem Maximum um die Mittagszeit und deutlichem Sommerüberschuss. Die Windenergie ist
wetterbestimmt, variiert über mehrtägige Wetterphasen und verzeichnet ein Winterhoch. Beide Profile
ergänzen einander komplementär, was für die Auswahl der Visualisierungen richtungsweisend ist.

\textbf{Residuallast und Dunkelflaute.} Der nach Abzug der erneuerbaren Einspeisung verbleibende Teil
der Last heißt Residuallast und muss von steuerbaren Kraftwerken, Speichern oder Importen gedeckt
werden. Fallen Wind und Sonne über Stunden bis Tage gleichzeitig aus, während die Last hoch ist
(\enquote{Dunkelflaute}), wird das gesamte konventionelle Reservoir gebraucht. Solche Lagen sind
selten und für die Systemauslegung dennoch prägend.

\textbf{Merit-Order und negative Preise.} Am Day-Ahead-Markt der Gebotszone DE-LU bestimmt das
Grenzkraftwerk den einheitlichen Börsenpreis. Erneuerbare Anlagen speisen mit Grenzkosten nahe null
ein und verschieben die Angebotskurve nach rechts. Ein hohes erneuerbares Angebot senkt entsprechend
den Preis. Übersteigt das Angebot die Nachfrage bei unflexibler Abregelung, fallen die Preise unter
null. Negative Börsenpreise bilden damit das reguläre Marktsignal für einen akuten Erzeugungsüberschuss.

\textbf{Grenzüberschreitender Handel.} Überschüsse fließen in das europäische Verbundnetz ab,
während Defizite durch Importe ausgeglichen werden. Der Nettohandelswert zeigt als dritter Baustein
neben Erzeugung und Preis, ob ein Überschuss im Inland verbleibt oder abfließt.

In der öffentlichen Debatte treten diese Zusammenhänge meist als Behauptungen auf
(\enquote{Erneuerbare machen Strom billig}, \enquote{im Winter liefern sie nichts}), werden aber
selten an Daten geprüft. Genau diese Prüfung soll die Anwendung ermöglichen.

\subsection{Zielgruppen}

\textbf{Z1 (Studierende und Lehrende energienaher Fächer, primär).} Als primäre Zielgruppe bearbeiten
sie praxisnahe Aufgabenstellungen und suchen konkrete Beispieltage für Seminare oder Abschlussarbeiten.
Fachbegriffe wie Merit-Order, Residuallast sowie die Einheiten GW und EUR/MWh sind ihnen geläufig,
der Umgang mit Rohdaten oder relationalen Datenbanken hingegen selten. Ihr Informationsbedarf besteht in
reproduzierbaren Fallbeispielen mit exakten Kennwerten.

\textbf{Z2 (Fachjournalismus, Energieagenturen und kommunale Versorger).} Diese Gruppe prüft und bewertet
energiewirtschaftliche Aussagen, häufig unter Zeitdruck und für eine breite Öffentlichkeit. Sie beherrscht
den Umgang mit statistischen Kennzahlen, besitzt jedoch wenig Erfahrung in komplexer Zeitreihenanalyse.
Ihr Ziel liegt im schnellen Auffinden belastbarer Einzelfälle samt Kontext, etwa dem teuersten Tag des
Jahres am 12.\ Dezember mit durchschnittlich $395{,}7$\,EUR/MWh.

\textbf{Z3 (Fachlich interessierte Prosumer).} Betreiber privater Photovoltaikanlagen, Batteriespeicher
oder Wärmepumpen mit dynamischen Stromtarifen kennen den eigenen Erzeugungsverlauf aus Monitoring-Apps.
Ihr Interesse gilt dem Abgleich dieser persönlichen Beobachtungen mit dem makroökonomischen Gesamtsystem.

Die drei Gruppen bringen etablierte Sehgewohnheiten mit, an die die Anwendung gezielt anschließt. Von
\emph{Energy-Charts} \cite{EnergyCharts} und Agora ist die Konvention \enquote{Erneuerbare in Grün- und Blautönen}
vertraut. Die Anwendung übernimmt Grün und Blau für die Erneuerbaren-Skala und reserviert ein warmes Rot für
den Preis, damit beide Größen nicht verwechselt werden. Kalender- und Aktivitätsheatmaps sind als
\enquote{eine Zelle = ein Zeitabschnitt, dunkler = mehr} allgemein bekannt. Darauf baut die Stunden-Heatmap auf.
Parallele Koordinaten setzen wir dagegen nicht voraus. Sie sind die einzige Technik, die erklärt werden muss,
weshalb über der Ansicht ein einzeiliger Hinweis steht und der Erstkontakt über Hover mit Tooltip erfolgt.
Preise werden in EUR/MWh erwartet und nicht in ct/kWh, Zeitangaben in UTC wie in den Quelldaten. Alle Zahlen
erscheinen mit deutschem Dezimalkomma.

\subsection{Überblick und Beiträge}

Die Grundlage der Anwendung bilden die EnergyCharts-Daten \cite{EnergyCharts} für Deutschland 2024,
abgerufen über den PostgREST-Spiegel der Universität Halle. Aus zwei Quelltabellen für Erzeugung und Preis
entsteht ein Datensatz mit $8\,690$ Stundenwerten und $364$ Tagesprofilen. Drei Techniken aus den vorgegebenen
Bereichen sind eng mit den Analyseaufgaben verzahnt. Eine Zeitreihen-Übersicht über die zwölf Monate dient als
Einstieg und Filter (Bereich Scatterplot und Zeitreihen), eine pixelorientierte Stunden-Heatmap fungiert als
Tag-Stunden-Matrix für alle $8\,690$ Stundenwerte (Bereich Pixel- und Icon-Techniken), und parallele
Koordinaten über sieben Tagesattribute decken den Bereich mehrdimensionaler Darstellungen ab. Die drei
Ansichten sind über einen gemeinsamen Interaktionszustand gekoppelt.

Diese Arbeit leistet fünf wesentliche Beiträge.

\begin{itemize}
  \item[\textbf{B1}] Eine reproduzierbare Datenaufbereitung (Kapitel~\ref{sec:daten}), die
        das in der Quelle unvollständige Lastfeld rechnerisch rekonstruiert und die $94$ fehlenden
        Stunden transparent als Datenlücke darstellt.
  \item[\textbf{B2}] Eine aus den Anwendungsaufgaben begründete Technikauswahl
        (Abschnitt~\ref{sec:praesentation}), bei der jede Technik gegen drei echte Alternativen
        abgewogen wird, die dasselbe mentale Modell aufbauen würden.
  \item[\textbf{B3}] Eine durchgängige Kopplung der drei Ansichten
        (Abschnitt~\ref{sec:interaktion}) über Mehrfach-Monatsfilter, geteilten Hover-Zustand,
        Mehrfachauswahl von Tagen und UND-verknüpftes Achsen-Brushing.
  \item[\textbf{B4}] Drei belegte Anwendungsfälle (Kapitel~\ref{sec:anwendungsfaelle}),
        darunter der Nachweis, dass die Häufigkeit negativer Börsenpreise vor allem dem
        Solaranteil folgt und $69{,}6\,\%$ dieser Stunden in das Fenster von 10 bis 15\,Uhr fallen.
  \item[\textbf{B5}] Eine offene Diskussion der Grenzen bei Lastrekonstruktion,
        Vorzeichenkonvention des Handelswertes und Datenlücken, damit die gezeigten Muster nicht
        überinterpretiert werden.
\end{itemize}

\abb{uebersicht}{1.0}{Gesamtansicht der Anwendung: Monatsfilter und Kennzahlen oben, darunter die
drei gekoppelten Visualisierungen, rechts das Detailpanel.}
{Screenshot der vollständigen Anwendung im Startzustand (kein Monatsfilter): Kopfzeile mit
Leitfrage, Monatsschaltflächen, sechs Kennzahlenkarten, Zeitreihen-Übersicht, Stunden-Heatmap,
parallele Koordinaten und rechts das Tagesdetailpanel.}

\section{Daten}\label{sec:daten}

\subsubsection*{Quelle und Anreicherung}

Die Daten stammen aus dem Bestand von \emph{Energy-Charts} des Fraunhofer ISE \cite{EnergyCharts},
der für die Lehrveranstaltung als PostgreSQL-Datenbank mit PostgREST-Schnittstelle
\cite{PostgREST} unter \url{https://dbs.informatik.uni-halle.de/sciencedata} bereitsteht. Genutzt
und miteinander verbunden werden \texttt{energycharts\_totalpower} (Erzeugungsleistung je
Energieträger in 15-Minuten-Auflösung, gefiltert auf 2024 über \texttt{unix\_seconds}) und
\texttt{energycharts\_price} (Börsenpreise, gefiltert auf \texttt{market\_id = "DE-LU"}).

Die Verbindung beider Quellen ist bereits eine inhaltliche Anreicherung: Erst die Kopplung von
physikalischer Erzeugung und ökonomischem Preis über denselben Zeitstempel macht F2 und F4
beantwortbar. Eine weitere Anreicherung mit Wetterdaten des DWD (Globalstrahlung,
Windgeschwindigkeit) haben wir geprüft und bewusst nicht umgesetzt: Sie würde die Erklärung von der
Wirkung (eingespeiste Leistung) auf die Ursache (Wetter) verschieben und damit eine zweite,
meteorologische Aufgabenklasse eröffnen, für die die drei gewählten Ansichten nicht ausgelegt sind.
Abschnitt~\ref{sec:ausblick} greift den Punkt als Erweiterung auf.

\subsubsection*{Technische Bereitstellung}

Der Zugriff auf die Quelle erfolgt einmalig über das Skript
\texttt{scripts/fetch\_energycharts.py}: Basic Authentication am Endpunkt \texttt{/token}, danach
seitenweiser Abruf über \texttt{/rpc/user\_table\_get} ($500$ Zeilen je Anfrage). Zwei
Besonderheiten der Schnittstelle waren wesentlich und sind im Skript berücksichtigt: Numerische
Bereichsbedingungen (Zeitstempel) funktionieren zuverlässig über das Feld \texttt{where\_} mit den
Operatoren \texttt{>=} und \texttt{<}; Textgleichheit (\texttt{country\_id}, \texttt{market\_id})
dagegen über das Feld \texttt{filters} als Objekt, z.\,B.\ \texttt{\{"market\_id": "DE-LU"\}}.

Ergebnis ist die Datei \texttt{data/energycharts\_de\_2024.json} ($3{,}4$\,MB, UTF-8). Sie wird von
der Anwendung zur Laufzeit per HTTP nachgeladen: \texttt{index.html} führt ein \texttt{fetch()} auf
denselben Ursprung aus und übergibt das Ergebnis als Flag an \texttt{Elm.Main.init}. Die Anwendung
setzt also keine lokal gespeicherten Daten voraus und funktioniert unverändert, sobald HTML,
JavaScript und JSON gemeinsam ausgeliefert werden (etwa über GitLab-Pages).

Die Trennung von Abholen und Ausliefern hat drei Gründe: Zugangsdaten dürfen nicht in eine
clientseitige Anwendung gelangen; die Abgabe wäre sonst von Tokengültigkeit und
Serververfügbarkeit abhängig; und die Rohdaten sind mit über $35\,000$ Viertelstundenzeilen und
mehreren Dutzend Spalten für eine interaktive Darstellung zu breit. Der Preis dafür ist der
Verzicht auf eine Live-Aktualisierung, den Abschnitt~\ref{sec:ausblick} diskutiert.

\subsubsection*{Datenvorverarbeitung}

\emph{1. Normierung der Einheiten.} Die Quellspalten tragen den Suffix \texttt{\_in\_gw}, führen die
deutschen Werte aber in MW-Skala (Solar-Mittagswerte um $43\,000$). Das Skript teilt daher durch
$1000$. Ohne diesen Schritt wären alle Achsen um drei Größenordnungen falsch beschriftet; der
Fehler fiel erst beim Vergleich mit bekannten Größenordnungen des deutschen Systems auf.

\emph{2. Aggregation auf Stunden.} Aus vier Viertelstundenwerten wird der arithmetische Mittelwert
gebildet. Der Börsenpreis ist im betrachteten Jahr ein Stundenprodukt, die Erzeugung liegt in
Viertelstunden vor. Die Stunde ist damit die feinste Auflösung, in der beide Größen gemeinsam
definiert sind. Eine Kopplung auf Viertelstundenebene würde den Preiswert dreimal künstlich
wiederholen und Scheingenauigkeit erzeugen. Der Mittelwert ist hier der richtige Aggregator, weil
Leistung über die Stunde integriert der Arbeit entspricht.

\emph{3. Rekonstruktion der Last.} Das Feld \texttt{load\_in\_gw} ist in der verwendeten
\texttt{totalpower}-Sicht für Deutschland durchgehend leer. Da die Last die Bezugsgröße jedes
Anteilswertes ist, wird sie rekonstruiert:
\[
  \mathrm{Last} \;=\; \mathrm{Solar} + \mathrm{Wind}_{\text{onshore}} + \mathrm{Wind}_{\text{offshore}}
                      + \mathrm{Residuallast}.
\]
Die Identität folgt aus der Definition der Residuallast als \enquote{Last abzüglich
dargebotsabhängiger Einspeisung}. Sie ist eine Näherung, weil Biomasse, Wasserkraft und Geothermie
nicht abgezogen werden. Diese Träger zählen in der Quelle nicht zur Residuallastdefinition. Die
Konsequenz wird in Schritt~5 sichtbar.

\emph{4. Abgeleitete Attribute.} Je Stunde entstehen die erneuerbare Summe (Solar, Wind on-/offshore,
Biomasse, Laufwasser, Speicherwasser, Geothermie), die fossile Summe (Braunkohle, Steinkohle, Öl,
Kokereigas, Erdgas) und der Erneuerbarenanteil als $100 \cdot \mathrm{EE}/\mathrm{Last}$. Je
Kalendertag folgen daraus Mittelwerte, das Lastmaximum, Solar- und Windanteil an der mittleren Last
sowie die Zahl der Stunden mit negativem Preis. Die Tagesebene ist kein Komfort, sondern die
Analyseeinheit für parallele Koordinaten und Detailpanel: Eine einzelne extreme Stunde sagt wenig,
ein Tagesprofil charakterisiert eine Wetterlage.

\emph{5. Werte über 100\,\%.} In $385$ der $8\,690$ Stunden liegt der Erneuerbarenanteil über
$100\,\%$, im Maximum bei $138{,}6\,\%$. Das ist kein Fehler, sondern Folge von
Überschusserzeugung, die exportiert oder eingespeichert wird, und der Näherung aus Schritt~3. Diese
Werte werden nicht auf $100\,\%$ gekappt. Kappen würde genau die Extremsituationen unsichtbar
machen, nach denen F2 sucht. Stattdessen läuft die Farbskala der Heatmap bis $120\,\%$ und sättigt
darüber (Abschnitt~\ref{sec:vis2}).

Die Vorverarbeitung liegt in Python und nicht in Elm. Das ist eine bewusste Abweichung von der
Empfehlung der Aufgabenstellung. Sie gehört zum Abruf aus einer authentifizierten Datenbank, den
die Anwendung aus Sicherheitsgründen nicht selbst ausführen soll. Alles, was nach dem Laden
geschieht, ist in Elm implementiert: Monatsfilterung, Monatsaggregation, Achsenskalierung,
Brush-Auswertung und Kennzahlen. Diese Schritte reagieren unmittelbar auf geänderte Rohdaten; ein
neuer Jahrgang in der JSON-Datei erfordert keine Codeänderung.

\subsubsection*{Umfang, Lücken und kritische Eignung}

\begin{table}[H]
\centering\small
\caption{Umfang des erzeugten Datensatzes}\label{tab:umfang}
\begin{tabular}{ll}
\toprule
Eigenschaft & Wert \\
\midrule
Land / Gebotszone & Deutschland / DE-LU \\
Zeitraum & 1.1.2024 bis 31.12.2024, Zeitstempel in UTC \\
Stundenwerte & $8\,690$ (von $8\,784$ möglichen) \\
Tagesprofile & $364$ (von $366$ möglichen) \\
Stunden mit Preisangabe & $8\,690$ (vollständig) \\
Stunden mit negativem Preis & $441$ ($5{,}1\,\%$) \\
Stunden mit EE-Anteil über $100\,\%$ & $385$ ($4{,}4\,\%$) \\
Mittlerer EE-Anteil / mittlerer Preis & $55{,}1\,\%$ / $78{,}6$\,EUR/MWh \\
Spannweite EE-Anteil (Stunde) & $12{,}0\,\%$ bis $138{,}6\,\%$ \\
Spannweite Preis (Stunde) & $-135{,}45$ bis $936{,}28$\,EUR/MWh \\
\bottomrule
\end{tabular}
\end{table}

Dem Datensatz fehlen $94$ Stunden, die in der Quelle nicht enthalten sind: zwei zusammenhängende
Ausfälle von je $47$ Stunden (23.05.\ 13:00 bis 25.05.\ 11:59 und 12.10.\ 13:00 bis 14.10.\ 11:59,
jeweils UTC). Der 24.05.\ und der 13.10.\ fehlen dadurch vollständig, vier weitere Tage sind nur mit
12 oder 13 Stunden besetzt und ihre Tagesmittel entsprechend schwächer gestützt. Die Lücken werden
nicht interpoliert: In einer pixelorientierten Darstellung wären interpolierte Werte visuell nicht
von Messwerten zu unterscheiden. Damit ginge genau die Eigenschaft verloren, die eine dichte
Darstellung wertvoll macht, nämlich dass jede Zelle für eine Beobachtung steht. Stattdessen zeichnet
die Heatmap zuerst ein vollständiges Raster neutraler grauer Zellen und legt die Datenzellen
darüber; wo Daten fehlen, bleibt Grau sichtbar und ist in der Legende als \enquote{keine Daten}
beschriftet.

Für die Fragestellungen sind die Daten gut geeignet: Die Preisreihe ist lückenlos, die zeitliche
Auflösung liegt deutlich unter der kürzesten interessierenden Struktur (dem Tagesgang), und
Erzeugung und Preis stammen aus demselben Marktgebiet. Drei Einschränkungen sind zu nennen. Erstens
ist die Last rekonstruiert; alle Anteilswerte sind Näherungen, deren systematische Richtung bekannt
ist. Sie fallen tendenziell zu hoch aus. Zweitens ist die Vorzeichenkonvention von
\texttt{cross\_border\_electricity\_trading\_in\_gw} in der Quelle nicht dokumentiert; der Wert wird
deshalb neutral als Nettohandelswert geführt. Die Daten stützen die Lesart \enquote{negativ =
Export}: Am Tag mit dem höchsten Erneuerbarenanteil (4.2., $95{,}9\,\%$) beträgt er $-13{,}1$\,GW, in
der Dunkelflautenstunde am 6.11.\ dagegen $+13{,}4$\,GW. Drittens zeigen die Daten Koinzidenz, nicht
Kausalität; die Anwendung ist ein Werkzeug zur Hypothesenbildung im Sinne von \cite{Keim2008}, kein
Erklärungsmodell.

\section{Visualisierungen}

Dieses Kapitel leitet die Anwendung her: zuerst die konkreten Aufgaben und die dafür nötigen
mentalen Modelle (\ref{sec:aufgaben}), daraus die Designanforderungen (\ref{sec:anforderungen}),
dann die drei Visualisierungen mit ihren Kodierungen und den verworfenen Alternativen
(\ref{sec:praesentation}) und schließlich die Interaktion, die sie zu einer Anwendung verbindet
(\ref{sec:interaktion}).

\subsection{Analyse der Anwendungsaufgaben}\label{sec:aufgaben}

Aus F1 bis F4 und den Bedürfnissen von Z1 bis Z3 ergeben sich sechs Aufgaben, bewusst als
Handlungen und nicht als Diagrammwünsche formuliert:

\begin{itemize}
  \item[\textbf{A1}] \emph{Orientieren.} Einen Einstiegszeitraum wählen, ohne vorher zu wissen, was
        ihn auszeichnet. (Voraussetzung aller weiteren Aufgaben)
  \item[\textbf{A2}] \emph{Rhythmen erkennen.} Feststellen, ob eine Schwankung dem Tagesgang, dem
        Wetterverlauf oder der Jahreszeit folgt. (F1, F3)
  \item[\textbf{A3}] \emph{Ausreißer lokalisieren.} Einzelne extreme Stunden in $8\,690$ Werten
        finden und ihren Zeitpunkt exakt ablesen. (F2)
  \item[\textbf{A4}] \emph{Ausdehnung prüfen.} Entscheiden, ob ein Extremwert eine Einzelstunde
        oder ein anhaltendes Muster ist. (F2)
  \item[\textbf{A5}] \emph{Tage vergleichen.} Mehrere Tage über alle relevanten Attribute
        gleichzeitig gegenüberstellen und ihre Ähnlichkeit beurteilen. (F3, F4)
  \item[\textbf{A6}] \emph{Gezielt auswählen.} Alle Tage finden, die zugleich mehrere
        Wertbedingungen erfüllen. (F2, F4)
\end{itemize}

\paragraph{Welche mentalen Modelle helfen dabei?} Die Aufgaben setzen voraus, dass die Nutzenden im
Kopf eine tragfähige Vorstellung des Systems aufbauen. Vier Modelle sind relevant, und jedes wird
durch eine der Ansichten gestützt.

\emph{M1: Die Zeit als zweidimensionales Raster.} Um A2 zu lösen, muss man Tagesgang und Jahresgang
als zwei unabhängige Achsen denken können und nicht als eine lange Linie. Erst wenn \enquote{jeder
Tag hat dieselben 24 Stunden} im Kopf steht, wird die Frage \enquote{tritt das immer zur selben
Tageszeit auf?} überhaupt stellbar. Für A2 und A4 ist das Modell notwendig: Ohne es bleibt eine
Häufung von Extremwerten eine unstrukturierte Menge von Zeitpunkten. Eine Darstellung, die die Zeit
in ein Tag\,$\times$\,Stunde-Raster faltet, vollzieht die Faltung bereits im Bild, statt sie dem
Betrachter aufzubürden.

\emph{M2: Die Deckungsbilanz.} Für A3 und A5 muss klar sein, dass der Erneuerbarenanteil ein
Verhältnis ist. Er kann steigen, weil die Erzeugung steigt oder weil die Last fällt. Ohne dieses
Modell werden Werte über $100\,\%$ als Fehler gelesen. Es entsteht, wenn Last, Solar, Wind und
Anteil gleichzeitig auf eigenen Achsen sichtbar sind.

\emph{M3: Der Preis als Folge der Knappheit.} A3 und A5 erfordern die Vorstellung, dass der Preis
kein weiteres Wettermerkmal ist, sondern aus der Deckungslücke folgt. Ob das zutrifft, ist selbst
eine empirische Frage. Die Anwendung muss den Zusammenhang deshalb sichtbar machen, statt ihn zu
unterstellen, und stellt Preis und Anteil gemeinsam dar, ohne sie zu verrechnen.

\emph{M4: Der Tag als Punkt in einem Merkmalsraum.} A5 und A6 sind ohne dieses Modell nicht lösbar:
\enquote{Ähnlichkeit zweier Tage} ist nur definiert, wenn ein Tag als Wertebündel gedacht wird und
Ähnlichkeit als Nähe in diesem Raum. Es ist das anspruchsvollste Modell und den Zielgruppen als
einziges unvertraut. Es braucht daher die stärkste visuelle Unterstützung: eine Darstellung, in der
ein Tag ein als Ganzes wahrnehmbares Objekt ist und die Achsen ihre ursprünglichen Einheiten
behalten.

Die Reihenfolge A1 $\rightarrow$ A2/A3 $\rightarrow$ A4/A5 $\rightarrow$ A6 entspricht dem Ablauf
\enquote{overview first, zoom and filter, details on demand} \cite{Shneiderman1996}.

\subsection{Anforderungen an die Visualisierungen}\label{sec:anforderungen}

\begin{itemize}
  \item[\textbf{R1}] \emph{Einstieg ohne Vorwissen.} Eine Ansicht muss das ganze Jahr auf einen
        Blick zeigen und als Filter dienen. (A1)
  \item[\textbf{R2}] \emph{Zwei Zeitskalen gleichzeitig.} Tagesgang und Jahresgang müssen in einer
        einzigen Darstellung sichtbar sein, ohne dass eine weggemittelt wird. (A2, M1)
  \item[\textbf{R3}] \emph{Vollständigkeit ohne Aggregation.} Alle $8\,690$ Stundenwerte müssen
        gleichzeitig darstellbar sein, damit seltene Ereignisse nicht durch Mittelung verschwinden
        ($441$ negative Preisstunden sind $5{,}1\,\%$ der Daten). (A3)
  \item[\textbf{R4}] \emph{Zusammenhängende Muster erkennbar.} Benachbarte Stunden müssen räumlich
        benachbart erscheinen, damit eine mehrstündige Phase eine Fläche und keine Punktwolke
        ergibt. (A4)
  \item[\textbf{R5}] \emph{Mehrere Attribute mit verschiedenen Einheiten.} Mindestens sieben
        Attribute (GW, \%, EUR/MWh, Anzahl) müssen für einen Tag gemeinsam lesbar sein, ohne auf
        eine gemeinsame Skala gezwungen zu werden. (A5, M2, M4)
  \item[\textbf{R6}] \emph{Wertbasierte Auswahl.} Wertebereiche auf mehreren Attributen müssen
        kombinierbar und die Treffermenge sichtbar sein. (A6)
  \item[\textbf{R7}] \emph{Kopplung.} Jede Beobachtung muss in mindestens einer anderen Ansicht
        überprüfbar sein; eine Auswahl darf nicht auf eine Ansicht beschränkt bleiben.
        (\cite{Roberts2007})
  \item[\textbf{R8}] \emph{Genaue Werte auf Abruf.} Da aus einer Farbe kein Zahlenwert ablesbar
        ist, muss jede Markierung ihre Werte auf Anfrage exakt ausgeben. (A3, A5)
  \item[\textbf{R9}] \emph{Ehrliche Kodierung.} Fehlende Daten dürfen nicht wie Nullwerte
        aussehen; Farbskalen müssen wahrnehmungsgerecht geordnet sein
        \cite{Borland2007,Harrower2003}; die wichtigste Größe gehört auf den genauesten Kanal
        \cite{Ware2012,Munzner2014}.
\end{itemize}

\begin{table}[H]
\centering\small
\caption{Zuordnung von Ansichten, Anforderungen und Aufgaben}\label{tab:anforderungen}
\begin{tabular}{lll}
\toprule
Ansicht & erfüllt vorrangig & unterstützt Aufgaben \\
\midrule
Zeitreihen-Übersicht (Vis 1) & R1, R7, R8 & A1, teilweise A2 \\
Stunden-Heatmap (Vis 2) & R2, R3, R4, R9 & A2, A3, A4 \\
Parallele Koordinaten (Vis 3) & R5, R6, R7 & A5, A6 \\
Detailpanel und Tooltips & R8 & A3, A5 \\
\bottomrule
\end{tabular}
\end{table}

\subsection{Präsentation der Visualisierungen}\label{sec:praesentation}

Die drei Techniken stammen aus drei verschiedenen der vorgegebenen Bereiche. Entscheidend ist aber
nicht diese formale Erfüllung, sondern dass jede diejenige Anforderung erfüllt, an der die beiden
anderen scheitern: Die Zeitreihe kann R3 nicht erfüllen, weil sie mittelt, die Heatmap nicht R5,
weil sie eine Größe zeigt, und die parallelen Koordinaten nicht R2, weil sie keine Zeitachse kennen.
Die Ansichten sind komplementär und nicht austauschbar. Bei jeder Technik werden Expressivität
(zeigt sie genau die Daten?) und Effektivität (nutzt sie für die wichtigste Größe den
wahrnehmungsstärksten Kanal?) im Sinne von \cite{Munzner2014,Ware2012} diskutiert.

\subsubsection{Visualisierung Eins: Zeitreihen-Übersicht der Monate}\label{sec:vis1}

\paragraph{Kodierung.} Die Ansicht zeigt zwölf Monatswerte. Der mittlere Erneuerbarenanteil wird
als Balkenhöhe kodiert (Kanal: Position auf gemeinsamer Achse bzw.\ Länge), der mittlere Preis als
Polylinie mit quadratischen Markern auf einer zweiten, rechts beschrifteten Achse in EUR/MWh.
Gitterlinien bei $0/25/50/75/100/125\,\%$ stützen das Ablesen. Farbe trägt keine Datenbedeutung,
sondern nur die Zuordnung zur Größe: Grün für den Anteil, Rot für den Preis, erklärt durch eine
Legende. Klick auf Balken oder Marker schaltet den Monat im Filter an bzw.\ aus, Hover liefert die
Monatswerte inklusive Solar- und Windanteil und Summe der negativen Preisstunden.

\paragraph{Warum diese Kodierung?} Der Erneuerbarenanteil ist die Leitgröße und liegt deshalb auf
dem genauesten Kanal (Position/Länge auf gemeinsamer Basislinie); der Preis als Nebengröße auf einer
Linie, die den Verlauf betont statt des exakten Wertes. Die doppelte Achse ist eine bewusst
eingegangene Schwäche: Der visuelle Abstand zwischen Balken und Linie hat keine Bedeutung, und die
scheinbare Korrelation hängt von der willkürlichen Skalierung der zweiten Achse ab. Sie ist
vertretbar, weil die Ansicht nicht zum Ablesen von Zusammenhängen dient, sondern als Navigations-
und Kontextebene für R1. Für Zusammenhänge sind Heatmap und parallele Koordinaten zuständig. Um aus
der Schwäche keine Fehlaussage werden zu lassen, ist die Preisachse explizit beschriftet und die
Interpretation im Bericht auf qualitative Aussagen beschränkt.

\paragraph{Alternativen.} \emph{Vollständige Stundenzeitreihe als Liniendiagramm.} Sie wäre
expressiver und erfüllte R3. Bei $8\,690$ Punkten auf einer Bildschirmbreite fallen jedoch rund
zwölf Werte auf ein Pixel; Überzeichnung macht Extremwerte ununterscheidbar von Rauschen, und der
Tagesgang erscheint als geschlossenes Band. Für A1 ist sie unbrauchbar. Ersatzlos gestrichen wird
sie deshalb nicht, sondern durch die Heatmap ersetzt, die dasselbe Datenvolumen ohne Überzeichnung
zeigt.

\emph{Horizon Graph.} Er löst genau dieses Auflösungsproblem: Durch Falten und Einfärben des
Wertebereichs erreicht er bei gleicher Höhe eine mehrfach höhere Datendichte und ist für lange
Zeitreihen Stand der Technik \cite{Aigner2011}. Verworfen, weil sein Lesen geübt sein muss (Bänder,
Vorzeichenfaltung) und er damit R1 verletzt, also den Einstieg ohne Vorwissen für Z2 und Z3. Für ein
reines Expertenpublikum wäre er die bessere Wahl.

\emph{Cycle Plot oder monatliche Boxplots.} Sie zeigten die innermonatliche Streuung, die der
Mittelwertbalken unterschlägt. Das ist ein realer Nachteil des gewählten Designs und wird in
Abschnitt~\ref{sec:anwendung1} sogar Gegenstand eines Anwendungsfalls. Verworfen, weil die Ansicht
zugleich Schaltfläche ist: Ein Balken ist eine große, eindeutige Klickfläche für den Monatsfilter,
eine Box mit Whiskern nicht. Die verlorene Streuungsinformation ist in der Heatmap vollständig
vorhanden, die Klickbarkeit wäre nicht ersetzbar.

\abb{zeitreihe}{0.92}{Zeitreihen-Übersicht: Balken zeigen den mittleren Erneuerbarenanteil je Monat,
die rote Linie den mittleren Preis auf der rechten Achse. Der Tooltip nennt die exakten Monatswerte.}
{Screenshot der Monatsübersicht mit Tooltip über einem Monatsbalken: zwölf grüne Balken, rote
Preislinie mit Markern, rechte Preisachse in EUR/MWh, Legende und Tooltip mit EE-Anteil, Preis,
Solar- und Windanteil sowie negativen Preisstunden.}

\subsubsection{Visualisierung Zwei: Pixelorientierte Stunden-Heatmap}\label{sec:vis2}

\paragraph{Kodierung.} Die Heatmap ist die dichteste Ansicht und setzt R2, R3 und R4 um. Jede Stunde
des Jahres ist eine Zelle: Die $x$-Achse läuft über die Tage des gefilterten Zeitraums, die
$y$-Achse über die $24$ Tagesstunden (0\,h oben, 23\,h unten), die Füllfarbe kodiert den
Erneuerbarenanteil. Ungefiltert sind das $366$ Spalten zu je $24$ Zellen; die Zellbreite wird
dynamisch als $\min(30,\,1500/\text{Tage})$ Pixel berechnet, sodass die Darstellung bei
Monatsfilterung von einer Pixel- in eine Kacheldarstellung übergeht, ohne dass sich die Semantik
ändert. Das ist eine Umsetzung des pixelorientierten Prinzips nach Keim \cite{Keim2000}: möglichst
viele Datenobjekte gleichzeitig, ein Objekt pro Pixel, mit einer Anordnung, die der Semantik der
Daten folgt.

\paragraph{Warum diese Anordnung?} Die Faltung der eindimensionalen Zeit in ein zweidimensionales
Raster ist der eigentliche Beitrag dieser Ansicht: Sie realisiert M1 direkt im Bild. Eine Struktur,
die vom Wetter kommt, erscheint als vertikaler Streifen über alle Stunden eines Tages; eine
Struktur, die vom Sonnenstand kommt, als horizontales Band in denselben Tagesstunden. Beide Fälle
sind ohne jede Rechnung unterscheidbar. Darin liegt die Antwort auf F1 und F3.

\paragraph{Farbe.} Die Skala ist sequentiell und interpoliert von einem dunklen Blauton (niedriger
Anteil) über Blaugrün zu hellem Grün (hoher Anteil) mit gemeinsam steigender Helligkeit. Sie folgt
damit der bekannten Konvention \enquote{Erneuerbare in Grün/Blau} und vermeidet bewusst eine
Regenbogenskala, deren nicht-monotone Helligkeit künstliche Grenzen erzeugt und die Ordnung der
Werte verschleiert \cite{Borland2007,Harrower2003}. Der Wertebereich wird bei $120\,\%$ begrenzt und
in der Legende als \enquote{$\geq 120\,\%$} ausgewiesen, sodass die $385$ Stunden über $100\,\%$
sichtbar bleiben, ohne den Kontrast des Hauptbereichs aufzubrauchen. Fehlende Stunden erhalten ein
neutrales, unbuntes Grau, das in der Farbskala nicht vorkommt und in der Legende eigens erklärt ist
(R9). Dass Farbe ein schwächerer Kanal als Position ist, wird bewusst in Kauf genommen und durch R8
aufgefangen: Der Tooltip liefert für jede Zelle Datum, Stunde, EE-Anteil, Preis, Last, Solar, Wind
und Nettohandel als Zahlen.

\paragraph{Alternativen.} \emph{Kalender-Heatmap nach van Wijk und van Selow} \cite{vanWijk1999}:
Sie ordnet Tage in einem Kalenderraster (Wochentage $\times$ Wochen) an und färbt jeden Tag nach
einer Kennzahl. Sie ist die naheliegendste Alternative und für Wochenrhythmen wie Werktag gegen
Wochenende überlegen. Verworfen, weil sie einen Tag auf einen Wert reduziert und damit R2 verletzt:
Der Tagesgang, die für Solar entscheidende Struktur, verschwindet. Ihre Stärke ist hier
zweitrangig, weil die Erzeugung im Gegensatz zur Last kaum Wochenstruktur hat.

\emph{Überlagerte Tageskurven (\enquote{Spaghettiplot}).} $366$ Kurven über $24$ Stunden zeigen den
Tagesgang exakt und erlauben das Ablesen von Werten. Verworfen wegen Überzeichnung: Im
Mittagsbereich ist die Dichte nicht mehr aufzulösen, und die für A3 zwingende Zuordnung einer Kurve
zu einem Datum ist ohne Interaktion unmöglich. Die Heatmap verzichtet auf exakte Werte (Farbe statt
Position) und gewinnt dafür die vollständige, überlagerungsfreie Zuordnung jeder Stunde zu ihrem
Ort im Bild.

\emph{Small Multiples je Monat.} Zwölf kleine Tag\,$\times$\,Stunde-Raster wären für den Vergleich
zwischen Monaten stärker, weil jeder Monat einen eigenen Rahmen bekommt. Verworfen, weil der
Übergang zwischen Monaten in getrennten Rahmen zerschnitten würde, etwa das kontinuierliche
Auseinanderlaufen des solaren Bandes von Februar bis Juni. Da die Anwendung den Monatsvergleich
bereits über Vis~1 und den Mehrfachfilter unterstützt, ist der durchgehende Zeitstrahl der größere
Gewinn.

\abb{heatmap}{0.98}{Stunden-Heatmap für das gesamte Jahr 2024: $x$ = Tag, $y$ = Stunde,
Farbe = Erneuerbarenanteil. Solares Mittagsband, windgetriebene vertikale Streifen und die grauen
Datenlücken im Mai und Oktober sind unmittelbar erkennbar.}
{Screenshot der Heatmap über alle zwölf Monate ohne Filter, mit Legende (0 bis 120\,\% plus Eintrag
für Datenlücken) und den Stundenbeschriftungen 0\,h, 6\,h, 12\,h, 18\,h, 23\,h an der linken Achse.}

\subsubsection{Visualisierung Drei: Parallele Koordinaten der Tagesprofile}\label{sec:vis3}

\paragraph{Kodierung.} Jeder gefilterte Tag ist eine Polylinie über sieben parallele, vertikale
Achsen: mittlere Last (GW), Solaranteil (\%), Windanteil (\%), Erneuerbarenanteil (\%), Nettohandel
(GW), mittlerer Preis (EUR/MWh) und Anzahl negativer Preisstunden. Jede Achse ist unabhängig auf
Minimum und Maximum im sichtbaren Datenbestand skaliert; beide Extremwerte stehen als Zahl an den
Achsenenden. Die Linien sind mit niedriger Deckkraft gezeichnet, sodass sich häufige Verläufe durch
Überlagerung von selbst verdichten. Hover hebt eine Linie rot hervor, Auswahl schwarz, vom Brushing
ausgeschlossene Tage werden fast vollständig ausgeblendet. Ist ein Achsenfilter aktiv, werden die
Treffer zusätzlich in kräftigem Grün gezeichnet: Ohne diese Hervorhebung bliebe die Treffermenge auf
derselben niedrigen Deckkraft wie das ungefilterte Bündel, und der Filter wäre zwar wirksam, aber
kaum sichtbar.

\paragraph{Warum diese Technik?} Parallele Koordinaten \cite{Inselberg1990,Heinrich2013} sind die
einzige der drei Techniken, die M4 aufbaut: Ein Tag wird zu einem einzigen Objekt, das man als
Ganzes verfolgen kann, und Ähnlichkeit zweier Tage wird als Ähnlichkeit zweier Linienverläufe
unmittelbar wahrnehmbar. Der entscheidende Vorteil gegenüber jeder Projektion ist, dass die Achsen
ihre Originaleinheiten behalten. Ein abgelesener Wert ist eine physikalische Größe und keine
Komponente. Genau das macht die Technik zur Grundlage für R6: Ein aufgezogener Bereich auf der
Preisachse ist eine Bedingung an den Preis und nicht an eine abstrakte Dimension.

\paragraph{Expressivität und Effektivität.} Expressiv sind parallele Koordinaten hier vollständig:
Alle sieben Attribute jedes Tages sind ohne Verlust dargestellt. Effektiv sind sie eingeschränkt.
Erstens hängt das Erkennen von Zusammenhängen von der Achsenreihenfolge ab. Nur benachbarte Achsen
zeigen ihre Korrelation als Bündelung bei gleichgerichtetem Verlauf oder als Kreuzung bei
gegenläufigem. Die Reihenfolge ist deshalb nicht beliebig, sondern folgt der Kausalkette des
Anwendungsgebiets: Last $\rightarrow$ Solar $\rightarrow$ Wind $\rightarrow$ EE $\rightarrow$
Handel $\rightarrow$ Preis $\rightarrow$ negative Stunden. Dadurch liegen die Paare nebeneinander,
deren Zusammenhang die Fragestellungen betreffen; das Paar (EE, Preis) zeigt das Auseinanderlaufen
der Linien, das der gemessenen Korrelation von $r = -0{,}80$ auf Tagesebene entspricht. Zweitens ist
bei $364$ Linien Überzeichnung unvermeidbar; sie wird durch niedrige Deckkraft in eine
Dichteinformation umgewandelt und durch Brushing beherrschbar gemacht.

\paragraph{Alternativen.} \emph{Streudiagramm-Matrix (SPLOM).} Sie zeigt jedes Attributpaar exakt
und ist für Korrelationen stärker. Verworfen, weil sieben Attribute $21$ Panels ergeben, jedes mit
$364$ Punkten: Ein einzelner Tag ist dann $21$-mal vorhanden und muss gedanklich wieder
zusammengesetzt werden. Genau diese Leistung soll M4 abnehmen.

\emph{Projektion (PCA oder MDS) mit Streudiagramm.} Sie löst das Überzeichnungsproblem elegant:
$364$ Tage als $364$ Punkte, Ähnlichkeit als Nähe, Cluster unmittelbar sichtbar. Verworfen, weil die
Achsen ihre Bedeutung verlieren. Auf die Frage \enquote{warum liegen diese Tage zusammen?} antwortet
eine Projektion nicht, und eine Bedingung im Sinne von R6 wie \enquote{zeige Tage mit negativem
Mittelpreis und hohem Solaranteil} ist in einem Komponentenraum nicht formulierbar.

\emph{Stern-Glyphen oder Radarplots je Tag.} Sie halten den Tag als Objekt zusammen und stützen M4.
Verworfen, weil der Vergleich von $364$ Glyphen ein serieller Suchvorgang ist, während parallele
Koordinaten alle Tage in einem gemeinsamen Bezugssystem übereinanderlegen. Hinzu kommt, dass Flächen
in Radarplots von der Achsenreihenfolge abhängen und systematisch fehlinterpretiert werden.

\abb{parallele-koordinaten}{0.98}{Parallele Koordinaten der Tagesprofile mit aktivem Brush auf der
Achse \enquote{Neg.\ Std.}: Nur Tage mit vielen negativen Preisstunden bleiben sichtbar, die
Statuszeile nennt die Trefferzahl.}
{Screenshot der parallelen Koordinaten über alle Tage, mit aufgezogenem Bereich im oberen Drittel der
rechten Achse, abgeblendeten übrigen Linien, Statuszeile \enquote{$n$ von 364 Tagen sichtbar} und
Schaltfläche zum Zurücksetzen der Achsenfilter.}

\subsection{Interaktion}\label{sec:interaktion}

Die drei Ansichten sind keine drei Diagramme, sondern eine Anwendung. Der gesamte
Interaktionszustand liegt in einem einzigen Elm-Model, aus dem alle Ansichten lesen; Konsistenz ist
dadurch nicht Konvention, sondern strukturell erzwungen. Vier Ebenen sind umgesetzt.

\paragraph{Ebene 1: Monatsfilter mit Mehrfachauswahl.} Monate lassen sich über eine
Schaltflächenleiste oder durch Klick auf einen Balken der Zeitreihe an- und abwählen; die Auswahl
ist eine Menge, sodass beliebige Kombinationen möglich sind (Juni bis August, oder Februar zusammen
mit November für einen Gegensatzvergleich). Sie bestimmt, welche Stunden in der Heatmap und welche
Tage in den parallelen Koordinaten und den Kennzahlenkarten erscheinen. \emph{Zweck:} A1 und
Reduktion der Datenmenge für A2 bis A6.

\paragraph{Ebene 2: Hover mit geteiltem Fokus und Tooltip.} Das Überfahren einer Heatmap-Zelle
oder einer Linie schreibt den Tag als \texttt{hoveredDay} in das Model. Beide Ansichten lesen diesen
Zustand, sodass die Hervorhebung immer gleichzeitig erscheint: Die Zellen des Tages erhalten einen
dunklen Rahmen, die Linie wird rot und in den Vordergrund gehoben, und das Detailpanel wechselt auf
diesen Tag. Zusätzlich erscheint am Mauszeiger ein Tooltip mit den exakten Werten. \emph{Zweck:} R7
und R8 ohne Zustandsänderung. Eine Vermutung lässt sich prüfen, ohne die bestehende Auswahl zu
verlieren.

\paragraph{Ebene 3: Mehrfachauswahl von Tagen.} Ein Klick auf Zelle oder Linie nimmt den Tag in
eine Liste auf, ein erneuter Klick entfernt ihn. Ausgewählte Tage bleiben in beiden Ansichten
dauerhaft schwarz hervorgehoben und erscheinen im Detailpanel als einzeln abwählbare Liste mit
EE-Anteil und Preis. Für den zuletzt berührten Tag zeigt das Panel sechs Kennzahlen und ein
Stundenprofil als kleines Balkendiagramm. Dessen Bezugslinien bei $50$ und $100\,\%$ und die
Stundenmarken alle sechs Stunden machen aus der reinen Verlaufsform eine ablesbare Größe. Die
$100$er-Linie ist dabei die inhaltlich wichtige, weil die Einspeisung oberhalb davon die Last
übersteigt.

Ab zwei ausgewählten Tagen tritt ein Vergleichsblock hinzu. Er besteht aus einer Tabelle mit allen
sieben Tagesattributen nebeneinander, einschließlich der Zahl negativer Preisstunden, die in der
Einzelansicht fehlt. Darunter stehen die Stundenprofile der ausgewählten Tage untereinander, sodass
sich die Tagesformen direkt gegenüberstehen (Solarbogen gegen flaches Windband). Beide sind nach
Datum sortiert und nicht nach Klickreihenfolge, damit derselbe Vergleich unabhängig von der
Bedienreihenfolge dasselbe Bild ergibt. \emph{Zweck:} A5, also der direkte Vergleich mehrerer Tage,
der mit einer Einfachauswahl nicht möglich ist.

\paragraph{Ebene 4: Achsen-Brushing.} Auf jeder der sieben Achsen lässt sich mit gedrückter
Maustaste ein Wertebereich aufziehen \cite{Becker1987}. Tage außerhalb werden abgeblendet und von
der Mausinteraktion ausgenommen. Brushes auf mehreren Achsen werden UND-verknüpft, sodass sich
zusammengesetzte Bedingungen formulieren lassen; eine Statuszeile nennt laufend die Trefferzahl, ein
Klick auf eine Achse löscht deren Filter, eine Schaltfläche alle.

Der Brush-Zustand wird dabei nicht in der Ansicht ausgewertet, sondern in \texttt{Main}. Von dort
speisen sich auch die Kennzahlenkarten, die sich damit auf die Treffermenge beziehen, und die
Heatmap, in der die Nichttreffer zurücktreten. Entscheidend ist, dass die Heatmap dafür abblendet
und nicht filtert. Alle Tage bleiben als Spalten stehen und die Zeitachse damit lückenlos.
Abgeblendete Tage sind von den grauen Zellen echter Datenlücken (R9) klar unterscheidbar. Ohne diese
Rückwirkung bliebe die stärkste Filteroperation der Anwendung auf eine einzige Ansicht beschränkt
und R7 wäre an dieser Stelle verletzt. \emph{Zweck:} R6, R7 und A6.

\begin{table}[H]
\centering\small
\caption{Kopplung der Ansichten: Jede Aktion wirkt über ihre Ursprungsansicht hinaus}\label{tab:kopplung}
\begin{tabular}{@{}p{3.9cm}p{2.9cm}p{3.6cm}p{3.6cm}@{}}
\toprule
Aktion & Zeitreihe & Heatmap & Parallele Koord. \\
\midrule
Monat wählen (Leiste oder Balken)
  & Balken markiert
  & zeigt nur diese Tage, Zellen werden breiter
  & zeigt nur diese Tage, Achsen reskalieren \\[0.3em]
Zelle überfahren (Heatmap)
  & unverändert
  & Tagesspalte gerahmt, Tooltip
  & Tageslinie rot hervorgehoben \\[0.3em]
Linie überfahren (Par.\ Koord.)
  & unverändert
  & Tagesspalte gerahmt
  & Linie rot, Tooltip mit 7 Werten \\[0.3em]
Zelle oder Linie klicken
  & unverändert
  & Tagesspalte dauerhaft markiert
  & Linie dauerhaft schwarz \\[0.3em]
Bereich auf Achse aufziehen
  & unverändert
  & Nichttreffer abgeblendet, Zeitachse bleibt vollständig
  & Treffer grün, Nichttreffer abgeblendet, Trefferzahl \\
\bottomrule
\end{tabular}
\end{table}

\paragraph{Was bewusst nicht umgesetzt wurde.} \emph{Umordnen der Achsen} wäre die wirksamste
Einzelverbesserung der parallelen Koordinaten \cite{Heinrich2013}, weil nur benachbarte Achsen auf
Korrelation prüfbar sind. Wir haben es zurückgestellt, weil die fachlich begründete feste
Reihenfolge die wichtigen Paare bereits benachbart macht und das Ziehen von Achsen mit dem Ziehen
von Brushes um dieselbe Geste konkurriert. \emph{Umschalten der Heatmap-Metrik} auf den Preis
bräuchte wegen dessen zweiseitiger Skala (negativ bis $936$\,EUR/MWh) eine divergierende Farbskala
mit Nullpunkt; zwei wechselnde Farbsemantiken in derselben Fläche sind eine bekannte Fehlerquelle
(siehe Abschnitt~\ref{sec:ausblick}). \emph{Ein Zeitraum-Brush in der Zeitreihe} wäre flexibler,
doch der Monat ist hier die natürliche Einheit, ein Klick billiger als ein gezogener Bereich, und
für tagesgenaues Eingrenzen steht die Heatmap zur Verfügung.

\abb{verknuepfung}{0.98}{Kopplung in Aktion: Monatsfilter Juli, der 7.\ Juli ist in der Heatmap
angeklickt und wird gleichzeitig in den parallelen Koordinaten hervorgehoben.}
{Screenshot mit Heatmap und parallelen Koordinaten untereinander, Monatsfilter Juli aktiv, die
Spalte des 7.\ Juli umrahmt, die zugehörige Linie hervorgehoben und ein Tooltip mit den
Stundenwerten.}

\section{Implementierung}\label{sec:implementierung}

\subsection*{Gesamtaufbau}

Die Anwendung ist eine \texttt{Browser.element}-Anwendung nach der Elm-Architektur
\cite{ElmArchitecture}: ein zentrales \texttt{Model}, ein \texttt{Msg}-Typ für alle Ereignisse, eine
reine \texttt{update}- und eine reine \texttt{view}-Funktion. Der Code ist in sechs Module
gegliedert (Tabelle~\ref{tab:module}) und folgt einer klaren Schichtung: \texttt{Data} kennt nur
Daten; die \texttt{Views}-Module kennen Daten und ihre eigene Darstellung, aber weder das
\texttt{Model} noch den \texttt{Msg}-Typ; \texttt{Main} kennt alles und verbindet es.

\begin{table}[H]
\centering\small
\caption{Modulstruktur der Anwendung (Stand Commit \texttt{\repocommit})}\label{tab:module}
\begin{tabular}{@{}llr@{}}
\toprule
Modul & Verantwortung & Zeilen \\
\midrule
\texttt{Main.elm} & Model, Msg, update, Komposition der Ansichten & 324 \\
\texttt{Data.elm} & Typen, JSON-Decoder, Filter- und Suchfunktionen, Formatierung & 342 \\
\texttt{Views/TimeSeries.elm} & Visualisierung Eins (Monatsübersicht) & 424 \\
\texttt{Views/Heatmap.elm} & Visualisierung Zwei (Stunden-Heatmap) & 277 \\
\texttt{Views/ParallelCoordinates.elm} & Visualisierung Drei (parallele Koordinaten, Brushing) & 461 \\
\texttt{Views/Layout.elm} & Filterleiste, Kennzahlen, Detailpanel, Vergleich, Tooltip, CSS & 475 \\
\midrule
\multicolumn{2}{@{}l}{Summe Elm} & 2303 \\
\texttt{scripts/fetch\_energycharts.py} & Datenabruf und Vorverarbeitung & 279 \\
\bottomrule
\end{tabular}
\end{table}

Die Schichtung wird durch ein durchgehendes Muster erzwungen: Jedes Visualisierungsmodul definiert
einen eigenen \texttt{Config msg}-Record, der den benötigten Zustandsausschnitt und die
auszulösenden Nachrichten als \emph{Funktionen} enthält. Die Module sind dadurch über \texttt{msg}
parametrisiert und kennen den konkreten \texttt{Msg}-Typ nicht:

\begin{lstlisting}[caption={Config-Record der Heatmap: Zustandsausschnitt und Callbacks},label={lst:config}]
type alias Config msg =
    { hoveredDay : Maybe String
    , selectedDays : List String
    , onHoverDay : String -> Tooltip -> msg
    , onMove : Float -> Float -> msg
    , onLeave : msg
    , onToggleDay : String -> msg
    }
\end{lstlisting}

Der Preis sind etwas längere Aufrufstellen in \texttt{Main.view}; der Gewinn ist, dass sich eine
Ansicht ohne Kenntnis der übrigen Anwendung austauschen oder wiederverwenden lässt und dass Hover-
und Klickverhalten in allen drei Ansichten garantiert identisch sind.

\subsection*{Die wichtigsten Datenstrukturen}

\paragraph{Das Model.} Der gesamte Interaktionszustand steht in einem Record; die Typentscheidungen
bilden jeweils eine Interaktionsanforderung ab.

\begin{lstlisting}[caption={Zentrales Model mit allen Interaktionszuständen},label={lst:model}]
type alias Model =
    { dataset : Result String Dataset   -- Decoder-Fehler explizit im Typ
    , selectedMonths : Set Int          -- Menge: Mehrfachauswahl von Monaten
    , hoveredDay : Maybe String         -- fluechtiger Fokus (Hover)
    , selectedDays : List String        -- dauerhafte Auswahl, Reihenfolge zaehlt
    , tooltip : Maybe Tooltip           -- Position und Inhalt am Mauszeiger
    , brushes : Dict Int ( Float, Float )  -- Achsenindex -> Wertanteil lo..hi
    , dragging : Maybe ParallelCoordinates.Drag  -- laufender Brush-Vorgang
    }
\end{lstlisting}

\texttt{Set Int} statt \texttt{Maybe Int} macht die Mehrfachauswahl von Monaten zu einer Eigenschaft
des Typs. \texttt{List String} statt \texttt{Set String} für die Tage ist Absicht: Die Reihenfolge
bestimmt, welcher Tag als zuletzt gewählter im Detailpanel erscheint. \texttt{brushes} speichert
\emph{normierte} Anteile $0..1$ statt absoluter Werte --- dadurch bleibt ein Brush gültig, wenn sich
durch einen Monatswechsel die Achsenskalierung ändert, und die Auswertung ist von der Dimension
unabhängig. \texttt{dragging} trennt den laufenden vom abgeschlossenen Brush, sodass während des
Ziehens live gefiltert werden kann, ohne den bestätigten Zustand zu überschreiben. Dass
\texttt{dataset} ein \texttt{Result} ist, zwingt \texttt{view}, den Fehlerfall zu behandeln --- ein
nicht dekodierbarer Datensatz erzeugt eine Meldung, keinen leeren Bildschirm.

\paragraph{Die Datentypen.} \texttt{Hourly} (14 Felder) und \texttt{Daily} (13 Felder) bilden die
beiden Aggregationsstufen ab. Diese Trennung ist die wichtigste Entwurfsentscheidung des
Datenmodells: Weil die Tagesaggregate vorliegen, muss keine Ansicht bei jeder Interaktion
gruppieren. \texttt{priceEurMwh : Maybe Float} macht die mögliche Abwesenheit eines Preises im Typ
sichtbar, sodass jede Ansicht sie behandeln muss.

\paragraph{Ein Detail des Decoders.} \texttt{Json.Decode} bietet Kombinatoren nur bis
\texttt{map8}; beide Records haben mehr Felder. Statt eine Zusatzbibliothek einzubinden, ist
\texttt{andMap} in drei Zeilen selbst definiert und der Decoder als Kette geschrieben:

\begin{lstlisting}[caption={andMap erweitert map8 auf beliebig viele Felder},label={lst:andmap}]
andMap : Decoder value -> Decoder (value -> result) -> Decoder result
andMap valueDecoder functionDecoder =
    Decode.map2 (\fn value -> fn value) functionDecoder valueDecoder
\end{lstlisting}

\subsection*{Funktion der drei Visualisierungen im Code}

Alle drei Ansichten erzeugen SVG mit \texttt{elm-community/typed-svg}. Der wesentliche Vorteil
gegenüber handgeschriebenem SVG ist, dass Attribute typisiert sind: Ein Längenwert ohne Einheit oder
ein falsch geschriebenes Attribut ist ein Übersetzungsfehler und kein stilles Nichtzeichnen.

\emph{Zeitreihe.} \texttt{monthly} faltet die Tagesliste zu zwölf \texttt{MonthSummary}-Records;
Monate ohne Daten werden über \texttt{filterMap} ausgelassen statt als Null gezeichnet. Die
Balkenhöhe skaliert linear auf $0..130\,\%$, die Preislinie auf das beobachtete Minimum und Maximum
--- die Preisachse ist damit datenabhängig und wird bei jedem Filterwechsel neu beschriftet.

\emph{Heatmap.} Die Zeichenfläche entsteht in drei Lagen: Stundenbeschriftungen, ein vollständiges
Raster grauer \enquote{Fehlt}-Zellen und darüber die Datenzellen. Diese Reihenfolge ist der Trick,
mit dem R9 ohne Sonderfallbehandlung erfüllt wird --- wo keine Datenzelle liegt, bleibt Grau stehen.
Die Farbfunktion bildet den Anteil auf einen Rampenparameter $t$ ab:

\begin{lstlisting}[caption={Farbrampe der Heatmap: Begrenzung bei 120\,\% statt Kappen der Daten},label={lst:color}]
renewableColor : Float -> Color.Color
renewableColor value =
    let
        t = clamp 0 1 (value / 120)
        r = 0.12 + (1 - t) * 0.72
        g = 0.25 + t * 0.47
        b = 0.42 + (1 - abs (t - 0.5) * 2) * 0.2
    in
    Color.rgb r g b
\end{lstlisting}

\emph{Parallele Koordinaten.} \texttt{dimensions} erzeugt die sieben Achsen als Records aus
Beschriftung, Zugriffsfunktion und beobachtetem Wertebereich; die Achsen sind damit datengetrieben,
eine achte Dimension wäre eine Zeile. \texttt{passesBrushes} prüft einen Tag mit \texttt{List.all}
gegen alle aktiven Brushes, was die UND-Verknüpfung direkt ausdrückt.

\subsection*{Funktion der Interaktion im Code}

Alle Ereignisse laufen über einen \texttt{Msg}-Typ mit zwölf Konstruktoren und werden in einer
einzigen \texttt{update}-Funktion behandelt. Die Kopplung entsteht dabei nicht durch Nachrichten
\emph{zwischen} Ansichten, sondern dadurch, dass alle aus demselben Zustand lesen --- das ist der
architektonische Kern der Anwendung.

Die aufwendigste Interaktion ist das Brushing, weil es Mauskoordinaten in Datenwerte übersetzen
muss. Der übliche Weg über \texttt{getBoundingClientRect} ist in Elm nur über Ports oder
\texttt{Browser.Dom}-Tasks erreichbar. Die Anwendung vermeidet das mit einer geometrischen
Entscheidung: Das SVG der parallelen Koordinaten wird in \emph{fester} Pixelgröße gerendert
($1280 \times 450$), sodass \texttt{offsetY} eines Mausereignisses direkt der SVG-$y$-Koordinate
entspricht. Damit reicht eine reine Funktion:

\begin{lstlisting}[caption={Umrechnung der Mausposition in einen normierten Achsenwert},label={lst:brush}]
yToFraction : Float -> Float
yToFraction offsetY =
    1 - clamp 0 1 ((offsetY - marginTop) / innerH)
\end{lstlisting}

Weil der Brush nicht nur die eigene Ansicht filtert, liegt seine \emph{Auswertung} nicht im
Views-Modul: \texttt{ParallelCoordinates.matchingDates} liefert die Menge der Datumswerte, die alle
gesetzten Bereiche passieren, und \texttt{Main} leitet daraus den Ausschnitt für Kennzahlen und
Heatmap ab. Der laufende Ziehvorgang geht dabei bereits ein, sodass live gefiltert wird. Das Modul
behält damit die Definition seiner Achsen --- nur dort ist bekannt, wie ein Pixelbereich in einen
Wertebereich übersetzt wird ---, gibt das Ergebnis aber als reine Funktion nach außen.

Damit \texttt{offsetY} browserübergreifend stimmt, beginnt die unsichtbare Trefferfläche jeder Achse
bei $y = 0$ und nicht erst am Achsenanfang --- andernfalls bezieht sich \texttt{offsetY} in manchen
Browsern auf die Oberkante des getroffenen Elements. Während eines Ziehvorgangs legt sich zusätzlich
ein transparentes Rechteck über die gesamte Zeichenfläche, das \texttt{mousemove}, \texttt{mouseup}
und \texttt{mouseleave} auffängt; ohne dieses Overlay bricht der Brush ab, sobald der Zeiger die
schmale Achsenspur verlässt.

\subsection*{Implementierungsaufwand}

\paragraph{Was leicht ging.} Der Aufbau der SVG-Diagramme --- Achsen, Gitter, Skalierungsfunktionen,
Pfadaufbau aus Punktlisten --- ließ sich weitgehend aus dem Code der Übungen übernehmen. Ebenso war
die Elm-Architektur durch die Übungen vertraut, sodass das Grundgerüst an einem Abend stand.
Monatsfilter und Hover-Hervorhebung sind jeweils wenige Zeilen, weil sie nur ein Feld im Model
setzen.

\paragraph{Was aufwendig war.} Erstens \emph{die Daten selbst}, nicht die Visualisierung: Dass
\texttt{load\_in\_gw} leer ist und die \enquote{\_in\_gw}-Spalten MW-Werte führen, war nur über
Plausibilitätsvergleiche zu erkennen; bis dahin waren alle Anteilswerte um Größenordnungen falsch.
Auch die Unterscheidung zwischen \texttt{where\_} und \texttt{filters} in der Schnittstelle war nur
durch systematisches Ausprobieren zu finden. Zweitens \emph{das Brushing} wegen des beschriebenen
Koordinatenproblems: Die erste Fassung arbeitete mit responsivem SVG und lieferte je nach
Fensterbreite verschobene Auswahlbereiche; die Lösung (feste Pixelgröße, Trefferfläche ab $y = 0$,
Overlay) entstand erst nach mehreren Fehlversuchen und ist im Code kommentiert, damit sie nicht
versehentlich zurückgebaut wird. Drittens \emph{die Decoder} wegen der \texttt{map8}-Grenze.
Viertens \emph{der Umgang mit Datenlücken}: Fehlende Stunden einfach wegzulassen führte dazu, dass
die Heatmap Tage stillschweigend zusammenschob und der Zeitstrahl falsch wurde --- erst die Lage aus
vorgezeichneten Grauzellen löste das strukturell.

\paragraph{Leistung.} Ungefiltert zeichnet die Anwendung rund $8\,690$ Datenzellen, $8\,784$
Hintergrundzellen und $364$ Polylinien als SVG-Elemente. Das ist im Browser flüssig, weil die
Aggregate vorberechnet sind und \texttt{update} bei jeder Interaktion nur ein Feld ändert. Für
mehrere Jahre in Stundenauflösung wäre die Grenze erreicht; dann wäre die Heatmap auf Canvas
umzustellen.

\paragraph{Fehlerbehandlung.} Die beiden wahrscheinlichsten Fehlerfälle sind getrennt behandelt:
Schlägt das \texttt{fetch()} fehl (Datei fehlt, kein Webserver), zeigt die HTML-Seite die
HTTP-Meldung; lässt sich die Datei laden, aber nicht dekodieren, zeigt Elm den Decoder-Fehler im
Klartext. Der Metadatenblock der JSON-Datei führt Land, Marktgebiet, Jahr, Quelle,
Erzeugungszeitpunkt und die Hinweise zu Normierung und Lastrekonstruktion mit.

\subsection*{Einsatz von KI-Werkzeugen}

Teile des Codes sind mit Unterstützung eines KI-Assistenten entstanden. Das Vorgehen war durchgehend
dasselbe: Wir haben Anwendungsaufgabe, gewünschte
Kodierung und Randbedingungen vorgegeben, das Ergebnis übersetzt, im Browser gegen die Daten geprüft
und gezielt nachgebessert. Selbst erstellt und nicht generiert sind die inhaltlichen Entscheidungen
dieses Berichts: die Auswahl der drei Techniken und ihre Begründung, die Achsenreihenfolge der
parallelen Koordinaten, die Behandlung der Werte über $100\,\%$ und der Datenlücken sowie die
Auswahl der Anwendungsfälle. Geprüft wurde generierter Code auf drei Wegen: Übersetzung mit
\texttt{elm make} (der Compiler schließt eine ganze Fehlerklasse aus), Vergleich der angezeigten
Werte mit unabhängig in Python berechneten Kennzahlen und manuelle Bedienung aller
Interaktionspfade.

\section{Anwendungsfälle}\label{sec:anwendungsfaelle}

\subsection{Anwendung Visualisierung Eins: Grenzen aggregierter Monatsmittel}\label{sec:anwendung1}

Die Monatsübersicht zeigt auf den ersten Blick den erwarteten saisonalen Verlauf: Die durchschnittlichen Börsenpreise erreichen im Winterhalbjahr ihre Höchststände (November mit $114{,}0$ und Dezember mit $108{,}2$\,EUR/MWh) und fallen im Frühjahr auf ein Minimum (April mit $62{,}4$\,EUR/MWh). Der Erneuerbarenanteil verhält sich dazu tendenziell gegenläufig. Bemerkenswert ist dabei vor allem die geringe Schwankungsbreite der monatlichen Durchschnittswerte: Sie verharren über das gesamte Jahr hinweg in einem engen Intervall zwischen $44{,}4\,\%$ im November und $60{,}8\,\%$ im April, was einer Differenz von lediglich 16 Prozentpunkten entspricht. Die zugrunde liegenden Stundenwerte streuen hingegen zwischen $12{,}0\,\%$ und $138{,}6\,\%$. Die wesentliche Dynamik des Systems findet folglich innerhalb der einzelnen Monate statt und wird durch reine Monatsmittelwerte verdeckt.

\begin{table}[H]
\centering\small
\caption{Monatswerte 2024, wie sie die Tooltips der Zeitreihen-Übersicht ausgeben}\label{tab:monate}
\begin{tabular}{@{}lrrrrrr@{}}
\toprule
Monat & Tage & Ø EE-Anteil & Ø Solaranteil & Ø Windanteil & Ø Preis & Neg.\ Std. \\
      &      & (\%)        & (\%)          & (\%)         & (EUR/MWh) & \\
\midrule
Januar    & 31 & 56{,}5 & 4{,}6  & 39{,}6 & 76{,}6  & 16 \\
Februar   & 29 & 58{,}5 & 7{,}2  & 38{,}8 & 61{,}4  & 4 \\
März      & 31 & 54{,}7 & 13{,}6 & 28{,}4 & 64{,}7  & 12 \\
April     & 30 & 60{,}8 & 19{,}0 & 29{,}3 & 62{,}4  & 50 \\
Mai       & 30 & 59{,}7 & 26{,}2 & 20{,}0 & 65{,}9  & 78 \\
Juni      & 30 & 58{,}2 & 26{,}3 & 19{,}4 & 72{,}9  & 64 \\
Juli      & 31 & 57{,}2 & 26{,}8 & 18{,}1 & 67{,}7  & 81 \\
August    & 31 & 53{,}7 & 25{,}6 & 16{,}7 & 82{,}0  & 68 \\
September & 30 & 56{,}9 & 19{,}2 & 25{,}8 & 78{,}1  & 40 \\
Oktober   & 30 & 47{,}7 & 11{,}0 & 23{,}7 & 89{,}7  & 9 \\
November  & 30 & 44{,}4 & 5{,}4  & 27{,}7 & 114{,}0 & 11 \\
Dezember  & 31 & 52{,}5 & 4{,}2  & 36{,}1 & 108{,}2 & 8 \\
\bottomrule
\end{tabular}
\end{table}

Beim Vergleich der einzelnen Spalten in Tabelle~\ref{tab:monate} fällt zudem auf, dass die Häufigkeit negativer Börsenpreise vor allem an die Photovoltaik gekoppelt ist und weniger der gesamten Erneuerbarenquote folgt. Während Februar und April mit $58{,}5\,\%$ und $60{,}8\,\%$ einen fast identischen Gesamterwerbsanteil aufweisen, verzeichnen sie 4 gegenüber 50 negativen Preisstunden bei Solaranteilen von $7{,}2\,\%$ gegenüber $19{,}0\,\%$. Der Januar verzeichnet trotz des Jahreshöchstwerts beim Windanteil ($39{,}6\,\%$) lediglich 16 negative Stunden. Ursächlich hierfür ist das Erzeugungsprofil: Windenergie speist kontinuierlicher über Tag und Nacht ein, während Solaranlagen ihre Leistung auf wenige Mittagsstunden bündeln und dort punktuell erhebliche Überdeckungen verursachen.

Für die Zielgruppen liefert dieser Befund wesentliche Erkenntnisse. Fachleute und politische Akteure (Z1 und Z2) erhalten ein klares Argument gegen die pauschale Annahme, ein hoher Gesamtanteil erneuerbarer Energien führe automatisch zu Preisverfall oder Marktverzerrungen. Für Betreiber flexibler Verbraucher (Z3) folgt daraus, dass dynamische Stromtarife im Frühjahr und Sommer signifikant mehr Niedrigpreisfenster bieten als in windstarken Wintermonaten. Methodisch unterstreicht die Beobachtung die Notwendigkeit nachgelagerter Sichten: Die Monatsübersicht identifiziert die Diskrepanz, deren Klärung jedoch eine stundenfeine Auflösung erfordert.

Die tabellarische Darstellung der zwölf Monatswerte erlaubt zwar das exakte Nachschlagen einzelner Zahlen, erfordert jedoch ein gezieltes Vergleichen der Wertepaare. Im Zeitreihendiagramm hingegen sticht die Gleichförmigkeit der Balkenhöhen bei gleichzeitig stark variierender Preislinie unmittelbar visuell hervor. Bei zwölf Beobachtungen bleibt der Mehraufwand einer Tabelle noch überschaubar. Der Nutzen der grafischen Abstraktion steigt jedoch mit zunehmender Anzahl an Perioden. Dies bestätigt die Rolle der Zeitreihe als übergeordnete Orientierungs- und Navigationsebene.

\subsection{Anwendung Visualisierung Zwei: Zeitliche Konzentration negativer Marktpreise}\label{sec:anwendung2}

In der Jahres-Heatmap treten zwei orthogonale Musterstrukturen hervor: Von März bis September erstreckt sich ein horizontales helles Band über die Mittagsstunden von etwa 9 bis 16 Uhr, das sich im Hochsommer ausdehnt und zum Herbst hin verengt. Dieses Band repräsentiert den deterministischen Sonnenverlauf. Quer dazu verlaufen schmalere, vertikale Streifen über alle 24 Stunden eines Tages, die über das gesamte Jahr verteilt sind und im Winter gehäuft auftreten. Sie spiegeln mehrtägige Windfronten wider. Beide Phänomene lassen sich visuell ohne rechnerische Vorverarbeitung direkt anhand ihrer Ausrichtung im Zeitraster unterscheiden.

Ein markanter Einzelfall lässt sich an der Zelle mit dem höchsten Erneuerbarenanteil des Jahres ablesen: Am 1.\ Mai 2024 um 12:00 Uhr erreichte der Anteil $138{,}6\,\%$ bei einem Börsenpreis von $-120{,}07$\,EUR/MWh. Beim Überfahren dieser Zelle blendet die Anwendung einen Tooltip ein, der die exakten Messwerte dieser Stunde ausgibt (Abbildung~\ref{fig:anwendungsfall-mai}). Zu dieser Stunde stand einer Gesamtlast von $43{,}0$\,GW eine solare Einspeisung von $43{,}5$\,GW, eine Winderzeugung von $9{,}8$\,GW sowie ein Nettoexport von $-12{,}1$\,GW gegenüber. Die Solarleistung allein übertraf somit bereits den gesamten inländischen Bedarf. Gleichzeitig hebt der dunkle Rahmen die gesamte Tagesspalte des 1.\ Mai hervor und verdeutlicht die hohe Erzeugung auch in den angrenzenden Stunden. Ein Klick auf die Zelle überträgt den Tag in das Detailpanel, das für den gesamten 1.\ Mai 8 negative Preisstunden bei einer mittleren Tagesdeckung von $94{,}8\,\%$ ausweist.

\abb{anwendungsfall-mai}{0.98}{Anwendungsfall zur Heatmap: Monatsfilter Mai, Tooltip auf der hellsten Zelle des Jahres (1.\ Mai, 12:00\,Uhr, $138{,}6\,\%$, $-120{,}07$\,EUR/MWh). Die grauen Spalten am 23.--25.\ Mai zeigen die Datenlücke der Quelle.}
{Screenshot der Heatmap mit Monatsfilter Mai und geöffnetem Tooltip über der Mittagszelle des 1.\ Mai, rechts im Bild die grauen Zellen der Datenlücke vom 23.\ bis 25.\ Mai.}

Eine systematische Auszählung aller 441 negativen Preisstunden im Jahresverlauf bestätigt das visuelle Bild: $69{,}6\,\%$ dieser Stunden fallen in das Zeitfenster zwischen 10 und 15 Uhr, mit einer deutlichen Spitze von 74 Stunden um 12 Uhr mittags. In den gesamten Nachtstunden von 0 bis 6 Uhr traten in Summe lediglich 53 negative Stunden auf. Negative Börsenpreise erweisen sich damit primär als Phänomen solarer Mittagsspitzen.

Diese Differenzierung besitzt für energieintensive Betriebe und Speicherbetreiber (Z3) hohe praktische Relevanz, da die Verlagerung von Lasten in das Mittagsfenster verlässlichere Einsparpotenziale bietet als der spekulative Betrieb in windstarken Nächten. Für energiewirtschaftliche Analysen (Z1) liefert die Heatmap den Beleg für die Auswirkung des Merit-Order-Effekts bei gleichzeitiger Einspeisung.

Ein alternatives Histogramm über die Tagesstunden könnte die zeitliche Häufung zwar quantitativ präziser darstellen, setzt jedoch die Vorab-Hypothese voraus, dass die Tageszeit die entscheidende Variable ist. Der Vorteil der zweidimensionalen Heatmap liegt im explorativen Entdecken: Die horizontale Struktur wird sichtbar, ohne dass gezielt danach gesucht werden musste. Eine klassische Kalenderansicht nach van Wijk und van Selow \cite{vanWijk1999} hätte den 1.\ Mai zwar als auffälligen Tag hervorgehoben, den stundenweisen Verlauf jedoch verschleiert. Ein fortlaufendes Liniendiagramm über alle $8\,690$ Stunden wiederum würde aufgrund starker Überzeichnung die Periodizität der Mittagsstunden verdecken.

\subsection{Anwendung Visualisierung Drei: Identifikation gegensätzlicher Systemzustände}\label{sec:anwendung3}

Werden in der Anwendung der 4.\ Februar und der 6.\ November ausgewählt, zeichnen die parallelen Koordinaten zwei Linienprofile, die auf der Achse des Erneuerbarenanteils die extremen Ränder des Jahres markieren und über alle übrigen Achsen gegenläufig verlaufen (Tabelle~\ref{tab:gegensatz} und Abbildung~\ref{fig:anwendungsfall-vergleich}).

\begin{table}[H]
\centering\small
\caption{Zwei Gegenprofile, wie sie Detailpanel und Tooltips ausgeben}\label{tab:gegensatz}
\begin{tabular}{@{}lrr@{}}
\toprule
Tagesattribut & 4.\ Februar 2024 & 6.\ November 2024 \\
\midrule
Ø EE-Anteil & $95{,}9\,\%$ & $15{,}6\,\%$ \\
Ø Solaranteil & $3{,}7\,\%$ & $3{,}6\,\%$ \\
Ø Windanteil & $78{,}9\,\%$ & $0{,}4\,\%$ \\
Ø Last & $54{,}6$\,GW & $63{,}0$\,GW \\
Ø Nettohandel & $-13{,}1$\,GW & $+12{,}4$\,GW \\
Ø Preis & $13{,}31$\,EUR/MWh & $231{,}22$\,EUR/MWh \\
Negative Preisstunden & $0$ & $0$ \\
\bottomrule
\end{tabular}
\end{table}

An beiden Tagen liegt der durchschnittliche Solaranteil bei nahezu identischen $3{,}7\,\%$ beziehungsweise $3{,}6\,\%$. Die Differenz von rund 80 Prozentpunkten bei der erneuerbaren Deckung ist somit vollständig auf die Windenergie zurückzuführen ($78{,}9\,\%$ gegenüber $0{,}4\,\%$). In den parallelen Koordinaten kreuzen sich die beiden Polygonzüge auf der Solarachse in einem gemeinsamen Punkt, während sie auf der benachbarten Windachse maximal divergieren. Dies widerlegt die Annahme, hohe Anteile erneuerbarer Energien seien auf die Sommermonate beschränkt: Hohe Deckungsgrade werden im deutschen Stromsystem über zwei unterschiedliche Pfade erreicht, die saisonal entgegengesetzt liegen.

\abb{anwendungsfall-vergleich}{0.98}{Anwendungsfall zu den parallelen Koordinaten: 4.\ Februar und 6.\ November gleichzeitig ausgewählt und in beiden Ansichten markiert, das Detailpanel stellt sie im Vergleichsblock gegenüber.}
{Screenshot ohne Monatsfilter mit zwei angeklickten Tagesspalten in der Heatmap, den zugehörigen schwarzen Linien in den parallelen Koordinaten und dem Detailpanel rechts, das beide Tage als Chips auflistet und darunter alle sieben Tagesattribute sowie beide Stundenprofile gegenüberstellt.}

Zusätzliche Systemmuster lassen sich durch interaktives Brushing aufdecken (Abbildung~\ref{fig:parallele-koordinaten}). Bei einer Selektion des oberen Wertebereichs auf der Achse der negativen Preisstunden isoliert die Anwendung wenige Tage mit extremer Häufung. An der Spitze steht der 7.\ Juli 2024 mit 16 negativen Preisstunden bei einem Erneuerbarenanteil von $74{,}5\,\%$ ($32{,}5\,\%$ Solar- und $25{,}4\,\%$ Windanteil). Dabei weisen die verbleibenden Linienbündel eine einheitliche Struktur auf. Sie erreichen auf der Solarachse ausnahmslos hohe Werte, während sie auf der Gesamtachse über einen weiten Bereich streuen. Dies bestätigt die Erkenntnis aus der Heatmap in einer unabhängigen Darstellung.

In tabellarischer Form ließe sich der direkte Vergleich zweier bekannter Tage zwar numerisch exakt durchführen. Die parallelen Koordinaten bieten demgegenüber zwei wesentliche Vorteile. Einerseits verorten sie beide Tage unmittelbar im Kontext der Gesamtdarstellung als Randpunkte der Verteilung, andererseits erlauben sie das Auffinden relevanter Tage ohne vorherige Kenntnis konkreter Datumsangaben. Eine manuelle Filterung in einer Tabelle über 364 Zeilen und sieben Spalten wäre zwar rechnerisch möglich, böte jedoch keine Rückmeldung über die relative Lage der Treffermenge im multiattributiven Raum. Eine lineare Dimensionsreduktion wie die Hauptkomponentenanalyse (PCA) würde die extremen Tage zwar ebenfalls als distanzierte Punkte abbilden, ließe jedoch keine unmittelbare Rückführung auf die treibende physikalische Einzeldimension zu.


\section{Verwandte Arbeiten}

Die Einordnung in den wissenschaftlichen Kontext erfolgt entlang zweier Richtungen. Dies umfasst einerseits Arbeiten mit vergleichbarem energiewirtschaftlichen Anwendungsfokus sowie andererseits methodische Ansätze zur visuellen Verknüpfung mehrdimensionaler und dichter Zeitreihendarstellungen.

\subsection*{Visuelle Analyse von Energie-Zeitreihen}

Eine direkte Parallele zur vorliegenden Arbeit bildet die Untersuchung von van Wijk und van Selow \cite{vanWijk1999}. Die Autoren aggregieren univariate Zeitreihen des Energieverbrauchs mittels Clusteranalyse zu mittleren Tagesprofilen und visualisieren die Zuordnung der Einzeltage in einem kalendarischen Raster. Beiden Ansätzen liegt die Prämisse zugrunde, dass der Tag die zentrale Analyseeinheit für periodische Energiesysteme darstellt und tages- sowie jahreszeitliche Dynamiken simultan sichtbar sein müssen. Dieser Gedanke motiviert im vorliegenden Projekt das Tag-Stunden-Raster der Heatmap sowie die tagesbasierten Polygonzüge im Koordinatensystem. Wesentliche Unterschiede bestehen im Abstraktionsgrad und in der Dimensionalität. Während van Wijk und van Selow Cluster-Repräsentanten univariater Daten automatisiert berechnen, zeigt die hiesige Anwendung alle 364 Tage im Detail und überlässt die Gruppenbildung der interaktiven Filterung. Zudem verarbeitet die vorliegende Arbeit sieben gekoppelte Dimensionen je Tag, was den Einsatz paralleler Koordinaten begründet.

Janetzko et al. \cite{Janetzko2014} analysieren Stromverbrauchsdaten von Gebäuden mithilfe unüberwachter Anomalieerkennung, wobei errechnete Unregelmäßigkeiten über der Zeit visualisiert werden. Ähnlich wie im vorliegenden Bericht sollen seltene Extremzustände in hochaufgelösten Zeitreihen unter Berücksichtigung des jeweiligen Tagestyps hervorgehoben werden. Die Ansätze unterscheiden sich jedoch in der Aufgabenverteilung. Während bei Janetzko et al. ein statistischer Algorithmus Auffälligkeiten markiert, definieren Anwender im hiesigen System über Brushing flexibel eigene Kriterien für relevante Zustände. Dies gewährleistet eine transparente Nachvollziehbarkeit für energiewirtschaftliche Bewertungen. Inhaltlich verknüpft das vorliegende Projekt zudem die physikalische Erzeugung direkt mit dem Marktpreis, während sich Janetzko et al. auf reine Verbrauchsdaten beschränken.

\subsection*{Dichte Darstellungen, mehrdimensionale Profile und Kopplung}

Die theoretische Basis der pixelbasierten Heatmap geht auf Keim \cite{Keim2000} zurück, der den Entwurf dichter Darstellungen als Optimierungsaufgabe zur Erhaltung semantischer Distanzen beschreibt. Das Prinzip, jedem Messwert ein visuelles Element ohne vorherige Glättung zuzuordnen, wird in der hiesigen Anwendung aufgegriffen. Während Keim primär Daten ohne inhärente Ordnungsstruktur betrachtet und raumfüllende Spiral- oder Peano-Layouts entwirft, gibt die Zeitstruktur des Stromsystems eine natürliche zweidimensionale Matrix aus Tag und Stunde vor. Die methodische Grenze dichter Farbraster, wonach quantitative Werte nur näherungsweise aus Farbstufen abgelesen werden können, wird durch bedarfsgerechte Tooltips und Kennzahlenblöcke kompensiert.

Heinrich und Weiskopf \cite{Heinrich2013} bieten einen umfassenden Überblick über Verfahren für parallele Koordinaten. In Übereinstimmung mit ihrer Taxonomie dient die Technik hier vor allem dem interaktiven Filtern und dem Vergleich individueller Profile gegenüber der Hintergrundverteilung. Auf weitergehende Visualisierungsformen wie Kantenbündelung oder Dichteschätzungen wurde zugunsten einer klaren Lesbarkeit verzichtet. Die Achsenreihenfolge ist fachlich begründet fixiert, um die unmittelbare Interpretation ohne Neukonfiguration zu unterstützen.

Das Kopplungskonzept stützt sich auf die Systematisierung koordinierter Mehrfachansichten nach Roberts \cite{Roberts2007}. Zentrale Mechanismen wie verknüpfte Selektion, Hervorhebung und Filterung basieren hierbei auf einem gemeinsamen Anwendungszustand, der durch die funktionale Elm-Architektur deklarativ verwaltet wird. Während Roberts allgemeine Frameworks für frei verschaltbare Ansichten beschreibt, setzt diese Arbeit auf drei fest aufeinander abgestimmte Sichten, die dediziert aus den fachlichen Fragestellungen abgeleitet sind.

Aus beiden Forschungslinien ergibt sich die Ausrichtung dieser Arbeit. Sie übernimmt den Tag als fundamentale Analyseeinheit für Energiedaten, verzichtet jedoch zugunsten der Nachvollziehbarkeit auf vorgeschaltete automatisierte Clusteringverfahren. Die Kombination aus zweidimensionaler dichter Zeitdarstellung und mehrdimensionalen Koordinaten ermöglicht eine hypothesenfreie Exploration, deren Mehrwert in der unmittelbaren empirischen Überprüfbarkeit der Systemzusammenhänge liegt.

\section{Zusammenfassung und Ausblick}\label{sec:ausblick}

\subsection*{Zusammenfassung und Mehrwert}

Die vorliegende Arbeit stellt eine interaktive Visual-Analytics-Anwendung in Elm vor, die $8\,690$ Stundenwerte und 364 Tagesprofile des deutschen Stromsystems aus dem Jahr 2024 erschließt. Der methodische Aufbau orientiert sich an sechs konkreten Anwendungsaufgaben und vier mentalen Modellen: Die aggregierte Zeitreihenansicht dient der übergeordneten Orientierung, die stundenfeine Heatmap visualisiert tages- und jahreszeitliche Periodizitäten, und die parallelen Koordinaten erlauben den simultanen Vergleich aller sieben Tagesattribute. Durch die Kopplung der Ansichten über einen zentralen Zustand entsteht eine zusammenhängende Analysemöglichkeit, bei der Interaktionen in einer Komponente unmittelbar auf die übrigen rückwirken.

Der praktische Nutzen für die Zielgruppen liegt in der Verifikation energiewirtschaftlicher Thesen anhand empirischer Daten. Dies verdeutlichen drei zentrale Ergebnisse:
Erstens belegen die Daten, dass monatliche Mittelwerte die Schwankungsbreite der erneuerbaren Deckung verdecken. Während die Monatsdurchschnitte über das Jahr lediglich zwischen $44{,}4\,\%$ und $60{,}8\,\%$ variieren, reicht die stündliche Spanne von $12{,}0\,\%$ bis $138{,}6\,\%$.
Zweitens traten $69{,}6\,\%$ aller negativen Preisstunden im Zeitfenster zwischen 10 und 15 Uhr auf, womit negative Börsenpreise eindeutig als spezifische Folge solarer Mittagsspitzen identifiziert werden.
Drittens markiert ein Wintertag (der 4.\ Februar mit $95{,}9\,\%$) den Jahreshöchstwert der Deckung, an dem der Solaranteil auf demselben niedrigen Niveau lag wie am Tag mit der geringsten Jahresdeckung ($15{,}6\,\%$ am 6.\ November). Die Differenz wird vollständig durch die Windeinspeisung getrieben. Alle Befunde lassen sich mit wenigen Interaktionsschritten reproduzieren und über die gekoppelten Ansichten nachvollziehen.

Ebenso wichtig für eine fundierte Interpretation sind die methodischen Grenzen: Die Anwendung konzentriert sich auf deskriptive statistische Koinzidenzen. Zudem basiert die Residuallast auf einer rechnerischen Rekonstruktion, und verbleibende Datenlücken der Primärquelle werden transparent als Leerstellen dargestellt.

\subsection*{Ausblick und Erweiterungsmöglichkeiten}

Für die Weiterentwicklung des Systems bieten sich Ansatzpunkte auf grafischer und datenseitiger Ebene:

Auf Ebene der Visualisierungskomponenten böte eine umschaltbare Farbskala in der Heatmap einen deutlichen Mehrwert. Würde dieselbe Tag-Stunden-Matrix wahlweise nach Börsenpreis oder grenzüberschreitendem Nettohandel eingefärbt, ließen sich Zusammenhänge zwischen Preisspitzen und Deckungslücken direkt räumlich vergleichen. Hierfür wäre eine divergierende Farbskala mit neutralem Nullpunkt erforderlich. Bei den parallelen Koordinaten würde eine dynamische Umordnung der Achsen die Erkennbarkeit korrelativer Muster zwischen nicht benachbarten Dimensionen verbessern. Für die Monatsübersicht wiederum böte sich die Integration kompakter Boxplots an, um die breite Stundenstreuung innerhalb der Monate bereits in der Übersichtsebene sichtbar zu machen.

Hinsichtlich der Datenbasis wäre eine Anreicherung mit meteorologischen Parametern des Deutschen Wetterdienstes sinnvoll, etwa Globalstrahlung, Windgeschwindigkeit und Umgebungstemperatur. Dies würde die Analysekette von den physikalischen Einflussgrößen über die Einspeisung bis hin zur Preisbildung schließen. Eine zusätzliche Aufschlüsselung der Stromflüsse nach Nachbarländern könnte über eine geografische Flusskarte visualisiert werden. Langfristig ermöglicht die Ausweitung auf mehrjährige Datenreihen die Beobachtung struktureller Veränderungen im Strommarkt, beispielsweise eine Ausweitung des Zeitfensters negativer Preise durch den weiteren Zubau an Photovoltaikanlagen. Technisch empfiehlt sich bei mehrjährigen Datensätzen die Umstellung der Heatmap von SVG auf HTML5-Canvas, um eine gleichbleibend hohe Interaktionsgeschwindigkeit zu gewährleisten.

\clearpage
\section*{Anhang: Git-Historie}
\addcontentsline{toc}{section}{Anhang: Git-Historie}

Das Projekt liegt unter \url{\repourl} (Branch: \texttt{\repobranch}, Bewertungsstand: Tag \texttt{\repotag}).

Die folgende Ausgabe entstand mit
\begin{verbatim}
git --no-pager log --all --graph --no-color --date=short \
    --pretty='format:%h%d (%s, %ad)'
\end{verbatim}

{\scriptsize\verbatiminput{git-historie.txt}}


\printbibliography

\end{document}
